\documentclass[11pt,a4paper]{article}
\usepackage[margin=1in]{geometry}
\usepackage{amsmath,amssymb,amsthm}
\usepackage{algorithm}
\usepackage{algpseudocode}
\usepackage{booktabs}
\usepackage{hyperref}

\theoremstyle{definition}
\newtheorem{definition}{Definition}[section]
\newtheorem{theorem}{Theorem}[section]
\newtheorem{lemma}[theorem]{Lemma}
\newtheorem{remark}{Remark}[section]

\title{\textbf{Rigorous Mathematical Specification, Channel Factorization, and Parameter Efficiency of Depthwise Separable Convolution}}
\author{Team Scaibu \\ \small Email: \texttt{jaydeep.9664920749@gmail.com} \\ \small Architecture and Core Engine Team}
\date{October 2026}

\begin{document}
\maketitle

\begin{abstract}
This document presents the formal mathematical derivation, spatial-channel factorization mechanics, and computational complexity bounds of Depthwise Separable Convolution. Standard convolutional layers perform joint spatial filtering and cross-channel linear combination simultaneously, incurring high computational and parameter costs. Depthwise separable convolution factors this joint operation into two decoupled stages: a lightweight depthwise spatial convolution per channel followed by a pointwise $1 \times 1$ inter-channel projection. We formulate the two-stage mapping, derive exact theoretical FLOP reduction ratios, detail step-by-step Big-O complexity, trace a worked numerical example, and document hardware memory bandwidth considerations.
\end{abstract}

\section{Notation and Mathematical Preliminaries}

Let $X \in \mathbb{R}^{C_{\text{in}} \times H_{\text{in}} \times W_{\text{in}}}$ denote the input activation tensor with $C_{\text{in}}$ channels and spatial grid dimensions $(H_{\text{in}}, W_{\text{in}})$.

\begin{itemize}
    \item $C_{\text{in}}, C_{\text{out}} \in \mathbb{N}_{\ge 1}$: Input and output channel counts.
    \item $K_h, K_w \in \mathbb{N}_{\ge 1}$: Depthwise spatial filter dimensions.
    \item $W_{\text{dw}} \in \mathbb{R}^{C_{\text{in}} \times K_h \times K_w}$: 3D depthwise filter kernel tensor.
    \item $W_{\text{pw}} \in \mathbb{R}^{C_{\text{out}} \times C_{\text{in}}}$: 2D pointwise ($1 \times 1$) projection matrix.
    \item $b \in \mathbb{R}^{C_{\text{out}}}$: Pointwise channel bias vector.
    \item $Z \in \mathbb{R}^{C_{\text{in}} \times H_{\text{out}} \times W_{\text{out}}}$: Intermediate depthwise spatial representation.
    \item $Y \in \mathbb{R}^{C_{\text{out}} \times H_{\text{out}} \times W_{\text{out}}}$: Final combined output feature tensor.
\end{itemize}

\begin{table}[h]
\centering
\caption{Summary of Mathematical Symbols for Depthwise Separable Convolution}
\begin{tabular}{lll}
\toprule
\textbf{Symbol} & \textbf{Type / Domain} & \textbf{Description} \\
\midrule
$X$ & $\mathbb{R}^{C_{\text{in}} \times H_{\text{in}} \times W_{\text{in}}}$ & Input feature map tensor \\
$W_{\text{dw}}$ & $\mathbb{R}^{C_{\text{in}} \times K_h \times K_w}$ & Depthwise spatial filter kernels \\
$W_{\text{pw}}$ & $\mathbb{R}^{C_{\text{out}} \times C_{\text{in}}}$ & Pointwise $1 \times 1$ linear projection weights \\
$Z$ & $\mathbb{R}^{C_{\text{in}} \times H_{\text{out}} \times W_{\text{out}}}$ & Intermediate spatially filtered tensor \\
$Y$ & $\mathbb{R}^{C_{\text{out}} \times H_{\text{out}} \times W_{\text{out}}}$ & Output convolved tensor \\
$s, p$ & $\mathbb{N}_{\ge 1}, \mathbb{N}_{\ge 0}$ & Spatial stride and padding parameters \\
\bottomrule
\end{tabular}
\end{table}

\section{Problem Definition and Core Concepts}

\subsection{Intuitive Meaning}
In a standard convolution, every output channel combines spatial filtering across all input channels simultaneously. Depthwise Separable Convolution hypothesizes that cross-channel correlations and spatial correlations can be decoupled. Spatial features are extracted within each channel individually, and channels are subsequently mixed using pure $1 \times 1$ linear projections, cutting compute and parameters by $8\times$ to $9\times$.

\subsection{Formal Mathematical Formulations}

\subsubsection{Stage 1: Depthwise Spatial Convolution}
For each channel $c_{\text{in}} \in \{0, \dots, C_{\text{in}}-1\}$ and coordinate $(h, w)$:
\begin{equation}
Z(c_{\text{in}}, h, w) = \sum_{i=0}^{K_h-1} \sum_{j=0}^{K_w-1} W_{\text{dw}}(c_{\text{in}}, i, j) X_{\text{pad}}(c_{\text{in}}, h \cdot s + i, w \cdot s + j)
\end{equation}

\subsubsection{Stage 2: Pointwise Channel Mixing ($1 \times 1$ Convolution)}
For each output channel $c_{\text{out}} \in \{0, \dots, C_{\text{out}}-1\}$:
\begin{equation}
Y(c_{\text{out}}, h, w) = b_{c_{\text{out}}} + \sum_{c_{\text{in}}=0}^{C_{\text{in}}-1} W_{\text{pw}}(c_{\text{out}}, c_{\text{in}}) Z(c_{\text{in}}, h, w)
\end{equation}

\section{Algorithm Specification}

\begin{algorithm}[H]
\caption{Depthwise Separable 2D Convolution Forward Pass}
\label{alg:depthwise_separable}
\begin{algorithmic}[1]
\Require Input $X \in \mathbb{R}^{C_{\text{in}} \times H_{\text{in}} \times W_{\text{in}}}$, depthwise weights $W_{\text{dw}}$, pointwise weights $W_{\text{pw}}$, stride $s$, padding $p$, bias $b$.
\Ensure Output $Y \in \mathbb{R}^{C_{\text{out}} \times H_{\text{out}} \times W_{\text{out}}}$, intermediate $Z$, dimensions $H_{\text{out}}, W_{\text{out}}$.
\State $H_{\text{out}} \gets \lfloor (H_{\text{in}} + 2p - K_h) / s \rfloor + 1$, $W_{\text{out}} \gets \lfloor (W_{\text{in}} + 2p - K_w) / s \rfloor + 1$
\State Allocate $Z \in \mathbb{R}^{C_{\text{in}} \times H_{\text{out}} \times W_{\text{out}}}$, $Y \in \mathbb{R}^{C_{\text{out}} \times H_{\text{out}} \times W_{\text{out}}}$
\For{$c_{\text{in}} \gets 0 \textbf{ to } C_{\text{in}} - 1$}
    \For{$h \gets 0 \textbf{ to } H_{\text{out}} - 1$}
        \For{$w \gets 0 \textbf{ to } W_{\text{out}} - 1$}
            \State $Z[c_{\text{in}}][h][w] \gets \sum_{i, j} W_{\text{dw}}[c_{\text{in}}][i][j] \cdot X_{\text{pad}}[c_{\text{in}}][h \cdot s + i][w \cdot s + j]$
        \EndFor
    \EndFor
\EndFor
\For{$c_{\text{out}} \gets 0 \textbf{ to } C_{\text{out}} - 1$}
    \For{$h \gets 0 \textbf{ to } H_{\text{out}} - 1$}
        \For{$w \gets 0 \textbf{ to } W_{\text{out}} - 1$}
            \State $Y[c_{\text{out}}][h][w] \gets b[c_{\text{out}}] + \sum_{c_{\text{in}}=0}^{C_{\text{in}}-1} W_{\text{pw}}[c_{\text{out}}][c_{\text{in}}] \cdot Z[c_{\text{in}}][h][w]$
        \EndFor
    \EndFor
\EndFor
\State \Return $\{Y, Z, H_{\text{out}}, W_{\text{out}}\}$
\end{algorithmic}
\end{algorithm}

\section{Theoretical Guarantees and Efficiency Ratio}

\begin{theorem}[Computational Complexity Reduction Ratio]
Let $H_{\text{out}}, W_{\text{out}}$ be the output spatial dimensions. The computational cost ratio $\rho$ of Depthwise Separable Convolution compared to standard 2D Convolution is:
\begin{equation}
\rho = \frac{\text{FLOPs}_{\text{DSConv}}}{\text{FLOPs}_{\text{StdConv}}} = \frac{1}{C_{\text{out}}} + \frac{1}{K_h K_w}
\end{equation}
\end{theorem}

\begin{proof}
The operations required by standard convolution are:
\begin{equation}
\text{FLOPs}_{\text{StdConv}} = H_{\text{out}} W_{\text{out}} \cdot C_{\text{in}} \cdot C_{\text{out}} \cdot K_h \cdot K_w
\end{equation}
The operations required by depthwise separable convolution are:
\begin{align}
\text{FLOPs}_{\text{DSConv}} &= \text{FLOPs}_{\text{dw}} + \text{FLOPs}_{\text{pw}} \\
&= H_{\text{out}} W_{\text{out}} \cdot C_{\text{in}} \cdot K_h \cdot K_w + H_{\text{out}} W_{\text{out}} \cdot C_{\text{in}} \cdot C_{\text{out}}
\end{align}
Dividing the two quantities:
\begin{equation}
\rho = \frac{H_{\text{out}} W_{\text{out}} C_{\text{in}} (K_h K_w + C_{\text{out}})}{H_{\text{out}} W_{\text{out}} C_{\text{in}} C_{\text{out}} K_h K_w} = \frac{K_h K_w + C_{\text{out}}}{C_{\text{out}} K_h K_w} = \frac{1}{C_{\text{out}}} + \frac{1}{K_h K_w}
\end{equation}
For $3 \times 3$ kernels ($K_h K_w = 9$) and $C_{\text{out}} \ge 64$, $\rho \approx \frac{1}{9} \approx 11.1\%$, representing an $8.9\times$ computational reduction.
\end{proof}

\subsection{Limitations}
\begin{itemize}
    \item \textbf{Arithmetic Intensity}: Depthwise layers have lower arithmetic intensity (FLOPs per byte transferred), requiring kernel fusion on GPU accelerators.
\end{itemize}

\section{Complexity Analysis}

\subsection{Time Complexity}
\begin{itemize}
    \item \textbf{Total Time Complexity}: $\Theta(C_{\text{in}} \cdot H_{\text{out}} \cdot W_{\text{out}} \cdot (K_h K_w + C_{\text{out}}))$ FLOPs.
\end{itemize}

\subsection{Space Complexity}
\begin{itemize}
    \item \textbf{Storage}: $\mathcal{O}((C_{\text{in}} + C_{\text{out}}) H_{\text{out}} W_{\text{out}})$ memory for intermediate and output tensors.
\end{itemize}

\section{Exactness and Error Analysis}

The factored matrix multiplications are algebraically exact.

\section{Worked Numerical Example}

Let $C_{\text{in}} = 2, H_{\text{in}} = 2, W_{\text{in}} = 2, C_{\text{out}} = 1, K = 2 \times 2, s = 1, p = 0$.
\begin{equation}
X_0 = \begin{bmatrix} 1.0 & 2.0 \\ 3.0 & 4.0 \end{bmatrix}, \quad X_1 = \begin{bmatrix} 2.0 & 2.0 \\ 2.0 & 2.0 \end{bmatrix}
\end{equation}
\begin{equation}
W_{\text{dw}, 0} = \begin{bmatrix} 1.0 & 1.0 \\ 1.0 & 1.0 \end{bmatrix}, \quad W_{\text{dw}, 1} = \begin{bmatrix} 0.5 & 0.5 \\ 0.5 & 0.5 \end{bmatrix}, \quad W_{\text{pw}} = [[1.0, 2.0]], \quad b = [0.0]
\end{equation}

\begin{enumerate}
    \item \textbf{Depthwise Pass}:
    \begin{align}
    Z(0, 0, 0) &= 1(1) + 1(2) + 1(3) + 1(4) = 10.0 \\
    Z(1, 0, 0) &= 0.5(2) + 0.5(2) + 0.5(2) + 0.5(2) = 4.0
    \end{align}
    \item \textbf{Pointwise Pass}:
    \begin{equation}
    Y(0, 0, 0) = 1.0 \cdot Z(0, 0, 0) + 2.0 \cdot Z(1, 0, 0) = 1.0(10.0) + 2.0(4.0) = 10.0 + 8.0 = 18.0
    \end{equation}
\end{enumerate}

\section{Numerical Considerations and Edge Cases}

\begin{itemize}
    \item \textbf{Bias Handling}: Bias is added only once during the pointwise mixing stage to avoid redundant bias offsets.
\end{itemize}

\begin{thebibliography}{9}
\bibitem{howard2017mobilenets} A.~G.~Howard et al., ``MobileNets: Efficient Convolutional Neural Networks for Mobile Vision Applications,'' \emph{arXiv preprint arXiv:1704.04861}, 2017.
\bibitem{chollet2017xception} F.~Chollet, ``Xception: Deep Learning with Depthwise Separable Convolutions,'' in \emph{IEEE Conference on Computer Vision and Pattern Recognition (CVPR)}, pp.~1251--1258, 2017.
\end{thebibliography}

\end{document}
